% !TEX root = main.tex
% Shared identities for the three separate step-learning settings in Appendix D.
\providecommand{\LRstep}{L_{\mathrm{step}}}
\providecommand{\LRout}{L_{\mathrm{out}}}
\providecommand{\LRsig}{\sigma}
\providecommand{\LRnorm}[1]{\left\lVert #1\right\rVert}
\providecommand{\LRinner}[2]{\left\langle #1,#2\right\rangle}
\theoremstyle{plain}
\newtheorem{lrlemma}[lemma]{Lemma}
\newtheorem{lrproposition}[proposition]{Proposition}
\newtheorem{lrtheorem}[theorem]{Theorem}
\newtheorem{lrremark}[remark]{Remark}
\crefname{lrlemma}{Lemma}{Lemmas}
\crefname{lrproposition}{Proposition}{Propositions}
\crefname{lrtheorem}{Theorem}{Theorems}
\crefname{lrremark}{Remark}{Remarks}

\subsection{Shared setup and gradient identities}
\label{app:weak-routing-initial}

All three settings use the step loss and innovation defined in
\cref{def:losses,sec:credit}. The target decomposition below
apply to any hidden representation with a linear readout. For the additional
positional identities and the full-prefix model, consider binary counting
$K=2$, with independent
$b_j\sim\operatorname{Bernoulli}(p)$, $0<p<1/2$, and uniform $s_0$.
The full-prefix result uses $p=1/4$; the supplied-routing result retains
its own general-$K$ assumptions. As in the main paper, the input
contains the prompt letters and the displayed running counts, and the
prediction at $i_t$ for $s_t$ sees $s_{t-1}$; the associated prompt-letter
site is $j_t$ and
\begin{equation}\label{eq:lr-constant-lag}
 i_t-j_t=n,\qquad t=2,\ldots,n.
\end{equation}
In the full-prefix setting, attention is normalized over the \emph{whole visible causal
prefix}, including query, separator, and displayed-state sites.  We write
$\chi=p(1-p)$, with every supervised transition weighted by $1/(n-1)$.

For a one-head Block 2 with hidden parameters $\phi$ and readout
$V$, let $h_t=h_\phi(x_t)$, $z_t=Vh_t$, and
\begin{equation}
 q_s=(1-p)e_s+pe_{s\oplus1},\qquad
 g_t=e_{s_t}-q_{s_{t-1}}=(b_t-p)(e_{s_{t-1}\oplus1}-e_{s_{t-1}}).
\end{equation}
The state $s_{t-1}$ is independent of the entire prompt-bit vector: this
follows from uniformity of $s_0$.  Consequently,
\begin{equation}\label{eq:lr-innovation-cov}
 \mathbb E[(b_t-p)(b_j-p)]=\chi\mathbf 1\{j=t\}.
\end{equation}
For fixed $p=1/4$, the right side is a positive constant, unlike the
outcome-correlation factor $\lambda^n$.

\begin{lrlemma}[Exact target decomposition and readout rule credit]
\label{lem:lr-decomposition}
For population step cross-entropy and arbitrary logits,
\begin{equation}\label{eq:lr-kl-decomp}
 \LRstep=H(p)+\mathbb E\operatorname{KL}
 (q_{s_{t-1}}\Vert\operatorname{softmax}(z_t))
 -\Phi_{\mathrm{rule}},\qquad
 \Phi_{\mathrm{rule}}=\mathbb E\LRinner{g_t}{z_t}.
\end{equation}
Writing $C(\phi)=\mathbb E[g_t h_t^\top]$, one has
$\nabla_V\Phi_{\mathrm{rule}}=C(\phi)$.  At an exact state-only predictor
$\operatorname{softmax}(z_t)=q_{s_{t-1}}$, one has
$\nabla_V\LRstep=-C(\phi)$.
\end{lrlemma}
\begin{proof}
The conditional expected label is $q_{s_{t-1}}$ and the centered
innovation is $g_t$.  Expanding cross-entropy as log-sum-exp minus the
target logit gives \eqref{eq:lr-kl-decomp}.  The readout derivatives follow
from $z_t=Vh_t$ and the usual softmax cross-entropy gradient.
\end{proof}

